\section{Conclusions}
\label{sec:conclusions}

The paper measures a protocol that makes the holdout number of a credit-risk validation trustworthy: a split declared per case, a rule that fixes the training window of a dated case, and three invariance gates, with a control against vacuous passes, that fail when a fitted parameter depends on rows outside the training information set. Under the protocol, the holdout AUC of a WoE scorecard on planted data has a bias of \nv{syn.w.none.bias}$\,\pm\,$\nv{syn.w.none.ci} and its 95\% interval covers the reference value in \nv{syn.w.none.cov}\% of \nv{syn.reps} repetitions; the same holds for the standardised logit, on 800 rows, on the card holders and for the product's own scorecard. When a step reads the holdout the cost depends on the step: nothing measurable for imputation, scaling and bin edges, \nv{syn.w.woe.bias} for the weight of evidence with holdout labels and \nv{syn.w.fit.bias} when every step reads it, with coverage falling to \nv{syn.w.fit.cov}\%; a calibration map fitted on the holdout adds \nv{syn.w.caltest.bias} and makes the calibration look better than it is. Five-fold cross-validation with the preprocessing outside the folds is the most optimistic procedure measured (\nv{syn.w.cvpre.bias} on the planted data and \nv{small.w.cvpre.bias} on the small sample).

At least one of the three gates fails for each of the \nv{h5.total} leaks planted in the fitting path in every one of \nv{e2.reps} draws, and none fails on the protocol. The gates are complementary (the first misses the calibration leaks, the second the input-only leaks, the third a choice among configurations in a single draw), and they detect a breach of a rule whatever its cost. For a choice among configurations on the holdout, repeated perturbation gives a detection of $1-K^{-d}$, which the simulation reproduces within \nv{e2.draws.zmax} standard errors.

A leak that exists only in development is a drop of discrimination out of time and not a drift of the feature: the test of the current AUC against the development value rejects in \nv{e3.s0.5.auc_test}\% of the repetitions at a shift of half a standard deviation, while the stability index of the leaked feature rejects in \nv{e3.s1.0.csi_test}\% and the single-feature screen in \nv{e3.s1.0.single_feature_flag}\% at a shift of one, which the second-order form of the index explains. On five cohorts of the loan data a random split inside the training window reproduces the level of the 5-year rate within \nv{e4.maxz} binomial standard errors (at most \nv{e4.holdse} points) and the declared out-of-time cohort shows gaps of up to \nv{e4.2006-07.sc.aoot_gap} points.

Of the nine predictions, \nv{h.held.word} held and \nv{h.notheld.word} did not. The inflation by the screen with holdout labels on the card holders narrows the claim about label-reading steps to the weight of evidence and the fit on everything; the out-of-time gap of at least two points holds for the scorecard and not for the gradient-boosting rung in two cohorts; and the contrast of cross-validation inside and outside the folds holds for the scorecard on every data set and not for the logit on the card holders (before the folds) and on the planted data (inside the folds). The predictions were written together with the outputs that decide them, and H6 and H8 were set with committed values in view. The next steps are the measurement of the same gates on a gradient-boosting and a neural challenger, the extension of the leaks to the features built from dated information, where the gates are blind, and the measurement of the gates on a case with panel structure, where the exchangeability that the first proposition needs does not hold.
